%% ============================================================
%%  Hochschul-Praktikumsdokumentation
%%  „Vom Wasserfall zum Sprint"
%%  Modul: MIS und Unternehmenssoftware
%% ============================================================
\documentclass[12pt, a4paper, twoside]{report}

%% ---- Encoding & Sprache ------------------------------------
\usepackage[utf8]{inputenc}
\usepackage[T1]{fontenc}
\usepackage[ngerman]{babel}
\usepackage{csquotes}

%% ---- Schrift -----------------------------------------------
\usepackage{lmodern}
\usepackage{microtype}

%% ---- Seitenlayout ------------------------------------------
\usepackage[
  top=2.5cm, bottom=2.5cm,
  left=3.0cm, right=2.5cm,
  headheight=15pt
]{geometry}
\usepackage{fancyhdr}
\pagestyle{fancy}
\fancyhf{}
\fancyhead[LE,RO]{\small\nouppercase{\leftmark}}
\fancyhead[RE,LO]{\small Vom Wasserfall zum Sprint}
\fancyfoot[C]{\thepage}
\renewcommand{\headrulewidth}{0.4pt}

%% ---- Grafik & Farbe ----------------------------------------
\usepackage{graphicx}
\usepackage[table,xcdraw]{xcolor}
\definecolor{hsPrimary}{RGB}{0,82,147}
\definecolor{hsAccent}{RGB}{220,230,242}
\definecolor{codeGray}{RGB}{245,245,245}
\definecolor{darkGreen}{RGB}{0,100,0}
\definecolor{darkRed}{RGB}{139,0,0}

%% ---- Tabellen ----------------------------------------------
\usepackage{booktabs}
\usepackage{tabularx}
\usepackage{multirow}
\usepackage{longtable}
\usepackage{array}
\newcolumntype{L}[1]{>{\raggedright\arraybackslash}p{#1}}
\newcolumntype{C}[1]{>{\centering\arraybackslash}p{#1}}
\newcolumntype{R}[1]{>{\raggedleft\arraybackslash}p{#1}}

%% ---- Code-Listings -----------------------------------------
\usepackage{listings}
\lstdefinestyle{javaStyle}{
  language=Java,
  backgroundcolor=\color{codeGray},
  basicstyle=\ttfamily\footnotesize,
  keywordstyle=\bfseries\color{hsPrimary},
  commentstyle=\itshape\color{darkGreen},
  stringstyle=\color{darkRed},
  numberstyle=\tiny\color{gray},
  numbers=left,
  stepnumber=1,
  numbersep=8pt,
  frame=single,
  framerule=0.4pt,
  rulecolor=\color{gray!50},
  breaklines=true,
  breakatwhitespace=false,
  showstringspaces=false,
  tabsize=4,
  captionpos=b,
  xleftmargin=10pt,
  xrightmargin=5pt,
  aboveskip=12pt,
  belowskip=8pt,
}
\lstset{style=javaStyle}

%% ---- TikZ / UML-ähnliche Darstellungen --------------------
\usepackage{tikz}
\usetikzlibrary{shapes.geometric, arrows.meta, positioning, fit, backgrounds, calc}

%% ---- Hyperlinks --------------------------------------------
\usepackage[
  colorlinks=true,
  linkcolor=hsPrimary,
  citecolor=hsPrimary,
  urlcolor=hsPrimary,
  pdftitle={Vom Wasserfall zum Sprint},
  pdfauthor={Praktikumsgruppe},
  pdfsubject={MIS und Unternehmenssoftware}
]{hyperref}

%% ---- Literatur ---------------------------------------------
\usepackage[
  backend=biber,
  style=authoryear,
  sorting=nyt,
  maxbibnames=3
]{biblatex}
\addbibresource{literatur.bib}

%% ---- Weitere Pakete ----------------------------------------
\usepackage{enumitem}
\usepackage{float}
\usepackage{caption}
\usepackage{subcaption}
\usepackage{amsmath}
\usepackage{amssymb}
\usepackage{pifont}
\usepackage{tcolorbox}
\tcbuselibrary{skins, breakable}
\usepackage{mdframed}
\usepackage{parskip}
\usepackage{setspace}
\usepackage{titlesec}
\usepackage{tocloft}

%% ---- Abstandseinstellungen ---------------------------------
\setstretch{1.3}
\setlength{\parskip}{6pt}

%% ---- Kapitelformatierung -----------------------------------
\titleformat{\chapter}[block]
  {\normalfont\LARGE\bfseries\color{hsPrimary}}
  {\thechapter.}{1em}{}
  [\vspace{2pt}\titlerule]
\titlespacing*{\chapter}{0pt}{30pt}{20pt}

\titleformat{\section}[block]
  {\normalfont\large\bfseries\color{hsPrimary!80!black}}
  {\thesection}{1em}{}
\titleformat{\subsection}[block]
  {\normalfont\normalsize\bfseries}
  {\thesubsection}{1em}{}

%% ---- Infoboxen --------------------------------------------
\newtcolorbox{infobox}[1][]{
  enhanced, breakable,
  colback=hsAccent, colframe=hsPrimary,
  fonttitle=\bfseries, title={#1},
  boxrule=0.8pt, arc=3pt,
  left=6pt, right=6pt, top=4pt, bottom=4pt,
}

\newtcolorbox{warningbox}{
  enhanced,
  colback=orange!10, colframe=orange!70!black,
  boxrule=0.8pt, arc=3pt,
  left=6pt, right=6pt, top=4pt, bottom=4pt,
}

\newtcolorbox{successbox}{
  enhanced,
  colback=green!8, colframe=darkGreen,
  boxrule=0.8pt, arc=3pt,
  left=6pt, right=6pt, top=4pt, bottom=4pt,
}

%% ---- Checkmarks / Kreuze ----------------------------------
\newcommand{\cmark}{\textcolor{darkGreen}{\ding{51}}}
\newcommand{\xmark}{\textcolor{darkRed}{\ding{55}}}

%% ---- Abkürzungen ------------------------------------------
\newcommand{\json}{\textsc{json}}
\newcommand{\uml}{\textsc{uml}}
\newcommand{\mis}{\textsc{mis}}
\newcommand{\oop}{\textsc{oop}}
\newcommand{\dod}{\textit{Definition of Done}}

%% ============================================================
\begin{document}

%% ---- Deckblatt ---------------------------------------------
\begin{titlepage}
  \begin{tikzpicture}[remember picture, overlay]
    \fill[hsPrimary] (current page.north west) rectangle
      ([yshift=-3.5cm] current page.north east);
    \fill[hsPrimary!20] (current page.south west) rectangle
      ([yshift=1.8cm] current page.south east);
  \end{tikzpicture}

  \vspace*{0.5cm}
  \begin{center}
    {\color{white}\Large\textbf{Hochschule}}\\[0.2cm]
    {\color{white}\large Fachbereich Wirtschaftsinformatik}\\[0.2cm]
    {\color{white}\normalsize Modul: MIS und Unternehmenssoftware}
  \end{center}

  \vspace{2.5cm}

  \begin{center}
    \begin{tcolorbox}[
      enhanced,
      colback=white, colframe=hsPrimary,
      boxrule=1.5pt, arc=5pt,
      width=0.88\textwidth,
      left=15pt, right=15pt, top=12pt, bottom=12pt,
    ]
      \centering
      {\huge\bfseries\color{hsPrimary} Vom Wasserfall zum Sprint}\\[0.5cm]
      {\large Praktikumsdokumentation}\\[0.4cm]
      \rule{0.6\textwidth}{0.4pt}\\[0.4cm]
      {\normalsize\itshape Entwicklungsmodelle im Vergleich:}\\[0.2cm]
      {\normalsize Wasserfallmodell und Scrum anhand eines}\\[0.2cm]
      {\normalsize konsolenbasierten Elektroniklagerverwaltungssystems}
    \end{tcolorbox}
  \end{center}

  \vspace{2.0cm}

  \begin{center}
    \begin{tabular}{ll}
      \textbf{Veranstaltung:}  & MIS und Unternehmenssoftware \\[4pt]
      \textbf{Semester:}       & Sommersemester 2025 \\[4pt]
      \textbf{Abgabedatum:}    & 08. Mai 2025 \\[4pt]
      \textbf{Betreuung:}      & Lehrveranstaltungsleitung \\[8pt]
      \hline\\[-4pt]
      \textbf{Gruppe:}         & Praktikumsgruppe 4 \\[4pt]
      \textbf{Mitglieder:}     & Denise Louise Doumbouya \\[2pt]
                               & Jonas Meier \\[2pt]
                               & Priya Nambiar \\[4pt]
      \textbf{Betreuer:}       & Prof.\ Dr.\ Andreas Hoffmann \\
    \end{tabular}
  \end{center}

  \vfill
  \begin{center}
    {\small\color{gray} Praktikumsdokumentation eingereicht im Rahmen des Bachelor-Studiengangs\\
    Wirtschaftsinformatik, Sommersemester 2025}
  \end{center}
\end{titlepage}

%% ---- Ehrenwörtliche Erklärung ------------------------------
\cleardoublepage
\thispagestyle{empty}
\vspace*{3cm}
\begin{center}
  \textbf{\large Ehrenwörtliche Erklärung}
\end{center}
\vspace{1cm}
Wir erklären hiermit ehrenwörtlich, dass wir die vorliegende Praktikumsdokumentation selbstständig
und ohne fremde Hilfe angefertigt haben. Alle Stellen, die wörtlich oder sinngemäß aus
veröffentlichten oder unveröffentlichten Quellen entnommen wurden, sind als solche kenntlich gemacht.

\vspace{2cm}
\begin{flushleft}
  Ort, Datum: \underline{\hspace{6cm}} \\[1.5cm]
  Unterschriften: \\[2cm]
  \begin{tabular}{p{5cm} p{5cm}}
    \rule{4.5cm}{0.4pt} & \rule{4.5cm}{0.4pt} \\
    Denise L.\ Doumbouya & Jonas Meier \\[1.5cm]
    \multicolumn{2}{l}{\rule{4.5cm}{0.4pt}} \\
    \multicolumn{2}{l}{Priya Nambiar}
  \end{tabular}
\end{flushleft}

%% ---- Inhaltsverzeichnis ------------------------------------
\cleardoublepage
\tableofcontents

%% ---- Abbildungsverzeichnis ---------------------------------
\cleardoublepage
\listoffigures
\addcontentsline{toc}{chapter}{Abbildungsverzeichnis}

%% ---- Tabellenverzeichnis -----------------------------------
\cleardoublepage
\listoftables
\addcontentsline{toc}{chapter}{Tabellenverzeichnis}

%% ============================================================
%%  KAPITEL 1: EINLEITUNG
%% ============================================================
\cleardoublepage
\chapter{Einleitung}
\label{chap:einleitung}

Softwareentwicklungsprojekte scheitern häufig nicht an einem Mangel an technischem Know-how,
sondern an ungeeigneten Entwicklungsprozessen. Die Standish-Group dokumentiert in ihren
regelmäßig erscheinenden Chaos-Reports, dass ein erheblicher Anteil von IT-Projekten die
vereinbarten Zeit-, Budget- und Qualitätsziele verfehlt~\autocite{standish2020}.
Dieser Befund macht die Wahl des richtigen Entwicklungsmodells zu einer zentralen
ingenieurwissenschaftlichen Entscheidung.

Im Rahmen des Praktikums \textit{MIS und Unternehmenssoftware} wurden zwei der wirkungsmächtigsten
Paradigmen der Softwareentwicklung gegenübergestellt: das klassische \textbf{Wasserfallmodell}
nach Royce~\autocite{royce1970} und das agile \textbf{Scrum-Framework}~\autocite{schwaber2020}.
Der gewählte Anwendungskontext ist ein \textit{konsolenbasiertes Elektroniklagerverwaltungssystem}
in Java, das die typischen Anforderungen eines mittelständischen Elektronikhändlers abbildet:
Bestandsführung, Mindestbestandsüberwachung, Datenpersistenz und nutzerfreundliche
Konsoleninteraktion.

\section{Motivation und Problemstellung}

Die Frage, welches Entwicklungsmodell für ein konkretes Projekt geeignet ist, lässt sich nicht
abstrakt beantworten. Sie erfordert die praktische Erfahrung beider Ansätze. Genau diese
Erfahrung war das didaktische Ziel dieses Praktikums: Das System wurde zunächst vollständig
nach dem Wasserfallmodell geplant und implementiert; anschließend wurde es in zwei Scrum-Sprints
iterativ erweitert und auf Basis echten Kundenfeedbacks überarbeitet.

Der entscheidende Erkenntnismoment trat im Sprint Review nach Sprint~1 ein: Der simulierte
Auftraggeber \textemdash{} repräsentiert durch die Lehrveranstaltungsleitung in der Rolle des
Product Owners \textemdash{} wünschte eine strukturelle Änderung des Ausgabeformats (\json{}
statt \textsc{csv}), englischsprachige Labels für die spätere internationale Nutzung sowie
eine klarere Klassenarchitektur. Diese Änderungswünsche hätten im Wasserfallmodell erheblichen
Mehraufwand bedeutet; im Scrum-Kontext wurden sie ohne Prozessbruch in Sprint~2 aufgenommen.

\section{Aufbau der Dokumentation}

Die vorliegende Dokumentation gliedert sich wie folgt: Kapitel~\ref{chap:projektbeschreibung}
und~\ref{chap:zielsetzung} beschreiben das Projektvorhaben und die Zielsetzung.
Kapitel~\ref{chap:anforderungen} enthält die Anforderungsanalyse.
Die Kapitel~\ref{chap:entwicklungsmodelle} bis~\ref{chap:scrum} erläutern die theoretischen
Grundlagen beider Entwicklungsmodelle. Die Scrum-Artefakte und der konkrete Projektverlauf
werden in den Kapiteln~\ref{chap:backlog} bis~\ref{chap:review} dokumentiert.
Die technische Umsetzung findet sich in den Kapiteln~\ref{chap:architektur}
und~\ref{chap:implementierung}. Qualitätssicherung und Testkonzept behandelt
Kapitel~\ref{chap:tests}. Den Abschluss bilden der Modellvergleich (Kapitel~\ref{chap:vergleich}),
die Reflexion (Kapitel~\ref{chap:reflexion}) und das Fazit (Kapitel~\ref{chap:fazit}).

%% ============================================================
%%  KAPITEL 2: PROJEKTBESCHREIBUNG
%% ============================================================
\chapter{Projektbeschreibung}
\label{chap:projektbeschreibung}

\section{Ausgangssituation und Aufgabenstellung}

Die fiktive Firma \textit{ElektroTech GmbH} betreibt ein mittelgroßes Elektroniklager und
verwaltet ihren Bestand bislang ausschließlich über Tabellenkalkulationen. Ziel des Praktikums
war es, für diese Firma ein einfaches, aber robustes Bestandsverwaltungssystem zu entwickeln,
das auf der Konsole läuft und keine externe Datenbankinfrastruktur voraussetzt.

Die Vorgabe der Lehrveranstaltung sah explizit vor, denselben fachlichen Anwendungsfall
zweimal zu durchlaufen: einmal plangetrieben nach dem Wasserfallmodell und einmal iterativ
mit Scrum. Dadurch sollte die Gruppe die Unterschiede beider Ansätze nicht nur theoretisch
verstehen, sondern am eigenen Entwicklungsprozess erleben.

\section{Systemüberblick}

Das System verwaltet sechs Produktkategorien des Elektronikbereichs:

\begin{table}[H]
  \centering
  \caption{Verwaltete Produktkategorien des Elektroniklagers}
  \label{tab:kategorien}
  \begin{tabularx}{0.75\textwidth}{C{0.5cm} L{3.5cm} L{6cm}}
    \toprule
    \textbf{\#} & \textbf{Kategorie} & \textbf{Repräsentatives Artikel-Beispiel} \\
    \midrule
    1 & Monitor     & Dell UltraSharp 27\textquotedbl{} 4K, 599,00\,€ \\
    2 & SSD         & Samsung 980 Pro 1TB, 89,99\,€ \\
    3 & Router      & ASUS RT-AX88U WiFi~6, 229,00\,€ \\
    4 & Tastatur    & Logitech MX Keys, 99,00\,€ \\
    5 & Laptop      & Lenovo ThinkPad X1 Carbon, 1.499,00\,€ \\
    6 & Grafikkarte & NVIDIA GeForce RTX~4070, 649,00\,€ \\
    \bottomrule
  \end{tabularx}
\end{table}

Jeder Artikel wird durch folgende Attribute beschrieben:

\begin{table}[H]
  \centering
  \caption{Datenmodell eines Lagerartikels}
  \label{tab:datenmodell}
  \begin{tabularx}{\textwidth}{L{2.5cm} L{2.5cm} L{6.5cm} C{2cm}}
    \toprule
    \textbf{Attribut} & \textbf{Typ} & \textbf{Beschreibung} & \textbf{Pflicht} \\
    \midrule
    \texttt{id}       & \texttt{int}    & Eindeutiger, auto-inkrementierter Bezeichner & \cmark \\
    \texttt{name}     & \texttt{String} & Artikelbezeichnung (englisches Label)         & \cmark \\
    \texttt{quantity} & \texttt{int}    & Aktueller Lagerbestand (Stückzahl)            & \cmark \\
    \texttt{minimum}  & \texttt{int}    & Mindestbestand; löst Warnung aus              & \cmark \\
    \texttt{price}    & \texttt{double} & Verkaufspreis in Euro                         & \cmark \\
    \bottomrule
  \end{tabularx}
\end{table}

\section{Technischer Rahmen}

\begin{itemize}[leftmargin=1.5em]
  \item \textbf{Programmiersprache:} Java~17 (LTS)
  \item \textbf{Build-System:} Maven~3.9
  \item \textbf{Versionskontrolle:} Git (lokales Repository)
  \item \textbf{Persistenz:} \json{}-Datei (\texttt{warehouse.json})
  \item \textbf{Externe Bibliotheken:} \texttt{com.google.code.gson} 2.10.1 (JSON-Serialisierung)
  \item \textbf{Entwicklungsumgebung:} IntelliJ IDEA Community Edition
  \item \textbf{Betriebssystem:} Windows~11 / macOS Sonoma (cross-platform)
  \item \textbf{Testframework:} JUnit~5
\end{itemize}

\section{Projektorganisation}

Die Gruppe bestand aus drei Studierenden mit folgenden Scrum-Rollen:

\begin{table}[H]
  \centering
  \caption{Rollenverteilung im Scrum-Team}
  \label{tab:rollen}
  \begin{tabularx}{0.85\textwidth}{L{4cm} L{3.5cm} L{5cm}}
    \toprule
    \textbf{Name} & \textbf{Rolle} & \textbf{Verantwortung} \\
    \midrule
    Denise L.\ Doumbouya & Scrum Master     & Prozessverantwortung, Sprint-Moderation \\
    Jonas Meier          & Entwickler       & Kernimplementierung, Klassenarchitektur \\
    Priya Nambiar        & Entwicklerin     & JSON-Export, Tests, Fehlerbehandlung \\
    \midrule
    Lehrveranstaltung    & Product Owner    & Anforderungen, Sprint Review, Feedback \\
    \bottomrule
  \end{tabularx}
\end{table}

%% ============================================================
%%  KAPITEL 3: ZIELSETZUNG
%% ============================================================
\chapter{Zielsetzung}
\label{chap:zielsetzung}

\section{Fachliche Ziele}

Das primäre fachliche Ziel war die Entwicklung eines funktionsfähigen
Elektroniklagerverwaltungssystems, das die Grundoperationen der Bestandsführung vollständig
abdeckt. Im Einzelnen wurden folgende Funktionsziele definiert:

\begin{enumerate}[leftmargin=1.5em, label=\textbf{F\arabic*.}]
  \item Das System ermöglicht das Anlegen neuer Artikel mit allen Pflichtattributen.
  \item Bestehende Artikel können über ein nummeriertes Konsolenmenü angezeigt werden.
  \item Der Bestand eines Artikels kann sowohl erhöht als auch verringert werden.
  \item Negative Lagerbestände werden systemseitig abgefangen und verhindert.
  \item Bei Unterschreitung des artikelspezifischen Mindestbestands wird eine
        Konsolenwarnung ausgegeben.
  \item Der aktuelle Lagerbestand kann in eine \json{}-Datei exportiert werden.
  \item Alle Benutzerinteraktionen erfolgen über ein englischsprachiges Konsolenmenü.
\end{enumerate}

\section{Methodische Ziele}

Neben den fachlichen Zielen verfolgte das Praktikum explizit methodische Lernziele:

\begin{enumerate}[leftmargin=1.5em, label=\textbf{M\arabic*.}]
  \item Die Studierenden erfahren den vollständigen Wasserfallzyklus (Anforderung \textrightarrow{}
        Design \textrightarrow{} Implementierung \textrightarrow{} Test \textrightarrow{} Übergabe)
        am eigenen Projekt.
  \item Die Studierenden erleben Scrum als gelebten Prozess mit echten Artefakten (Product Backlog,
        Sprint Backlog, Increment).
  \item Die Gruppe reflektiert, an welchen Stellen das Wasserfallmodell zu
        \textit{Entwicklung am Kunden vorbei} führt.
  \item Die Studierenden bewerten eigenständig, welches Modell für welchen Projektkontext geeignet
        ist.
\end{enumerate}

\section{Qualitative Ziele}

\begin{infobox}[Qualitätskriterien laut Lehrveranstaltungsvorgabe]
  Sauberkeit des Codes (Clean Code nach Robert C.\ Martin), Vollständigkeit der Fehlerbehandlung,
  Nachvollziehbarkeit der Commit-Historie, Detailtiefe der Reflexion sowie die kritische
  Auseinandersetzung mit beiden Entwicklungsmodellen wurden als Bewertungskriterien festgelegt.
\end{infobox}

%% ============================================================
%%  KAPITEL 4: ANFORDERUNGSANALYSE
%% ============================================================
\chapter{Anforderungsanalyse}
\label{chap:anforderungen}

\section{Methodik der Anforderungserhebung}

Die Anforderungserhebung verlief in beiden Modellen unterschiedlich. Im Wasserfallmodell
wurden die Anforderungen in einer einzigen Sitzung vollständig erhoben und in einem
Pflichtenheft fixiert. Im Scrum-Kontext entstanden die Anforderungen iterativ: Eine erste
Version des Product Backlogs wurde gemeinsam mit dem Product Owner erstellt und nach dem
Sprint Review von Sprint~1 auf Basis des Feedbacks verfeinert.

\section{Funktionale Anforderungen}

\begin{longtable}{L{1cm} L{4.5cm} L{6cm} C{1.5cm} C{1.5cm}}
  \caption{Vollständige Liste der funktionalen Anforderungen}
  \label{tab:fanforderungen} \\
  \toprule
  \textbf{ID} & \textbf{Anforderung} & \textbf{Akzeptanzkriterium} & \textbf{Prio.} & \textbf{Sprint} \\
  \midrule
  \endfirsthead
  \multicolumn{5}{c}{\tablename~\thetable{} (Fortsetzung)} \\
  \toprule
  \textbf{ID} & \textbf{Anforderung} & \textbf{Akzeptanzkriterium} & \textbf{Prio.} & \textbf{Sprint} \\
  \midrule
  \endhead
  \midrule
  \multicolumn{5}{r}{\textit{Fortsetzung nächste Seite}} \\
  \endfoot
  \bottomrule
  \endlastfoot
  FA-01 & Artikel anlegen & Artikel erscheint in der Liste mit korrekten Attributen & Hoch & S1 \\
  FA-02 & Artikel anzeigen & Alle Artikel werden tabellarisch in der Konsole ausgegeben & Hoch & S1 \\
  FA-03 & Bestand erhöhen & Eingabe einer positiven Zahl erhöht \texttt{quantity} korrekt & Hoch & S1 \\
  FA-04 & Bestand verringern & Eingabe einer positiven Zahl verringert \texttt{quantity} & Hoch & S1 \\
  FA-05 & Negativbestand verhindern & System gibt Fehlermeldung aus; Bestand bleibt unverändert & Hoch & S1 \\
  FA-06 & Mindestbestandswarnung & Bei \texttt{quantity < minimum} erscheint Warnung in der Konsole & Mittel & S1 \\
  FA-07 & JSON-Export & Datei \texttt{warehouse.json} wird korrekt erzeugt & Hoch & S2 \\
  FA-08 & Englische Labels & Alle Menüpunkte und Ausgaben auf Englisch & Mittel & S2 \\
  FA-09 & Konsolenmenü & Nummernbasiertes Menü mit Eingabevalidierung & Hoch & S1 \\
  FA-10 & OOP-Struktur & Trennung von Modell, Logik und Präsentation & Mittel & S2 \\
\end{longtable}

\section{Nicht-funktionale Anforderungen}

\begin{table}[H]
  \centering
  \caption{Nicht-funktionale Anforderungen}
  \label{tab:nfanforderungen}
  \begin{tabularx}{\textwidth}{L{1cm} L{3.5cm} L{8.5cm}}
    \toprule
    \textbf{ID} & \textbf{Kategorie} & \textbf{Anforderung} \\
    \midrule
    NF-01 & Portabilität   & Lauffähig unter Windows und macOS ohne Konfigurationsänderungen \\
    NF-02 & Wartbarkeit    & Klare Paketstruktur; jede Klasse hat genau eine Verantwortlichkeit \\
    NF-03 & Robustheit     & Fehlerhafte Konsoleneingaben führen nicht zum Programmabsturz \\
    NF-04 & Lesbarkeit     & Sprechende Bezeichner; keine Magic Numbers \\
    NF-05 & Erweiterbarkeit & Neue Artikelkategorien können ohne Architekturänderung hinzugefügt werden \\
    \bottomrule
  \end{tabularx}
\end{table}

\section{Abgrenzung des Systems}

Das System umfasst bewusst keine persistente Datenbankanbindung, keine grafische
Benutzeroberfläche und keine Netzwerkfunktionalität. Diese Abgrenzung ist praxismotiviert:
Die Anforderungen eines realen Kleinstlagers in der Frühphase der Digitalisierung rechtfertigen
keine Datenbankinfrastruktur. Der \json{}-Export bildet eine solide Brücke zu späteren
Erweiterungsstufen.

%% ============================================================
%%  KAPITEL 5: ENTWICKLUNGSMODELLE
%% ============================================================
\chapter{Entwicklungsmodelle im Überblick}
\label{chap:entwicklungsmodelle}

Software Engineering kennt eine Vielzahl von Prozessmodellen, die sich in ihrem Umgang mit
Veränderung, Unsicherheit und Kundeninteraktion fundamental unterscheiden. Für das vorliegende
Praktikum wurden zwei prototypische Vertreter ausgewählt: das sequenzielle \textbf{Wasserfallmodell}
als paradigmatisches Beispiel plangetriebener Entwicklung und \textbf{Scrum} als prominentester
Vertreter agiler Rahmenwerke.

Die Gegenüberstellung ist in der Praxis relevant, weil viele Unternehmen aus klassischen
Entwicklungsstrukturen kommen und den Wechsel in agile Prozesse vollziehen müssen. Das Praktikum
simuliert genau diesen Übergang \textendash{} vom Wasserfall zum Sprint.

%% ============================================================
%%  KAPITEL 6: WASSERFALLMODELL
%% ============================================================
\chapter{Wasserfallmodell}
\label{chap:wasserfall}

\section{Theoretische Grundlagen}

Das Wasserfallmodell wurde 1970 von Winston Royce in seinem wegweisenden Paper
\textit{Managing the Development of Large Software Systems} beschrieben, obwohl Royce es
selbst als unzulänglich für reale Projekte bezeichnete~\autocite{royce1970}. In der
industriellen Praxis wurde es dennoch zum dominierenden Entwicklungsparadigma der 1970er
bis 1990er Jahre.

Das Modell gliedert die Softwareentwicklung in streng sequenzielle Phasen:

\begin{figure}[H]
  \centering
  \begin{tikzpicture}[
    node distance = 0.4cm,
    phase/.style = {
      rectangle, rounded corners=4pt,
      minimum width=7cm, minimum height=0.9cm,
      text centered, font=\small\bfseries,
      fill=hsPrimary, text=white, draw=hsPrimary!80!black
    },
    arrow/.style = {-{Stealth[length=6pt]}, thick, color=hsPrimary!60}
  ]
    \node[phase] (req)   {1. Anforderungsanalyse (Requirements)};
    \node[phase, below=of req, fill=hsPrimary!80] (des) {2. Systemdesign (Design)};
    \node[phase, below=of des, fill=hsPrimary!65] (imp) {3. Implementierung (Implementation)};
    \node[phase, below=of imp, fill=hsPrimary!50] (tes) {4. Test \& Integration (Testing)};
    \node[phase, below=of tes, fill=hsPrimary!35, text=black] (dep) {5. Deployment \& Wartung};

    \draw[arrow] (req.south) -- (des.north);
    \draw[arrow] (des.south) -- (imp.north);
    \draw[arrow] (imp.south) -- (tes.north);
    \draw[arrow] (tes.south) -- (dep.north);

    \node[right=0.5cm of req, font=\tiny, align=left, color=gray] {Pflichtenheft, Use Cases};
    \node[right=0.5cm of des, font=\tiny, align=left, color=gray] {Architektur, UML};
    \node[right=0.5cm of imp, font=\tiny, align=left, color=gray] {Coding, Unit Tests};
    \node[right=0.5cm of tes, font=\tiny, align=left, color=gray] {Integrationstests, UAT};
    \node[right=0.5cm of dep, font=\tiny, align=left, color=gray] {Release, Dokumentation};
  \end{tikzpicture}
  \caption{Phasenmodell des klassischen Wasserfallmodells}
  \label{fig:wasserfall}
\end{figure}

\section{Anwendung im Praktikum}

Im ersten Teil des Praktikums wurde das Lager\-verwaltungs\-system vollständig nach dem
Wasserfallmodell entwickelt. Dies bedeutete konkret:

\begin{description}[leftmargin=2em, style=nextline]
  \item[Phase 1 \textendash{} Anforderungsanalyse (Tag 1--2):] Alle Anforderungen wurden in
    einem gemeinsamen Sitzungsprotokoll festgehalten. Das Ergebnis war ein vollständiges,
    aber statisches Anforderungsdokument. Kein Raum für spätere Änderungen war vorgesehen.

  \item[Phase 2 \textendash{} Design (Tag 3):] \uml{}-Klassendiagramm, Paketstruktur und
    Datenmodell wurden festgelegt. Zu diesem Zeitpunkt war das Ausgabeformat als
    \textsc{csv} spezifiziert \textendash{} eine Entscheidung, die sich später als
    Fehlgriff herausstellen sollte.

  \item[Phase 3 \textendash{} Implementierung (Tag 4--7):] Die Klassen \texttt{Item},
    \texttt{Warehouse} und \texttt{Main} wurden implementiert. Der Code folgte strikt
    dem im Design festgelegten Plan.

  \item[Phase 4 \textendash{} Test (Tag 8):] Manuelle Funktionstests wurden durchgeführt
    und in einer Testprotokoll-Tabelle dokumentiert.

  \item[Phase 5 \textendash{} Übergabe (Tag 9):] Das fertige System wurde dem Product Owner
    präsentiert.
\end{description}

\section{Probleme im Wasserfallprojekt}
\label{sec:wasserfall-probleme}

Die Präsentation des Wasserfall-Ergebnisses offenbarte strukturelle Schwächen:

\begin{warningbox}
  \textbf{Kritischer Befund:} Das \textsc{csv}-Format entsprach nicht den Vorstellungen des
  Auftraggebers. Die deutschen Menü-Labels waren für den international geplanten Einsatz
  ungeeignet. Die \texttt{Main}-Klasse enthielt mehr als 250 Zeilen und verletzte das
  \textit{Single Responsibility Principle}. Diese Mängel waren \textbf{nicht in letzter Minute
  behebbar}, da das Modell keine eingebaute Feedback-Schleife vorsieht.
\end{warningbox}

Diese Erfahrung illustriert das klassische Wasserfallproblem: Kundenwünsche, die sich erst
beim Anblick eines lauffähigen Systems konkretisieren, können im sequenziellen Prozess nicht
mehr kosteneffizient integriert werden.

%% ============================================================
%%  KAPITEL 7: AGILE ENTWICKLUNG
%% ============================================================
\chapter{Agile Entwicklung}
\label{chap:agil}

\section{Das Agile Manifest}

Im Jahr 2001 formulierten 17 erfahrene Softwareentwickler das \textit{Agile Manifesto}~\autocite{beck2001}.
Es proklamiert vier Grundwerte:

\begin{tcolorbox}[
  enhanced, breakable,
  colback=hsAccent!60, colframe=hsPrimary,
  title={\textbf{Die vier Werte des Agilen Manifests}},
  fonttitle=\bfseries, boxrule=1pt, arc=4pt
]
  \begin{enumerate}[label=\arabic*., leftmargin=1.5em]
    \item \textbf{Individuen und Interaktionen} über Prozesse und Werkzeuge
    \item \textbf{Funktionierende Software} über umfassende Dokumentation
    \item \textbf{Zusammenarbeit mit dem Kunden} über Vertragsverhandlung
    \item \textbf{Reagieren auf Veränderung} über das Befolgen eines Plans
  \end{enumerate}
  \vspace{6pt}
  \textit{Wobei die rechten Punkte durchaus Wert haben, die linken aber höher geschätzt werden.}
\end{tcolorbox}

\section{Agile Werte im Praktikumskontext}
\label{chap:agilewerte}

Die vier Manifest-Werte materialisieren sich im vorliegenden Praktikum wie folgt:

\paragraph{Individuen und Interaktionen:}
Anstelle starrer Rollentrennungen arbeiteten alle Gruppenmitglieder eng zusammen. Tägliche
Kurzabstimmungen (Daily Standup) \textendash{} auch wenn nur per Nachrichtengruppe durchgeführt
\textendash{} ersetzten formale Statusberichte.

\paragraph{Funktionierende Software:}
Das Ziel jedes Sprints war ein ausführbares, testbares Inkrement, kein Dokument. Nach Sprint~1
konnte der Product Owner das System tatsächlich starten und testen.

\paragraph{Zusammenarbeit mit dem Kunden:}
Das Sprint Review nach Sprint~1 ermöglichte echte Interaktion mit dem Product Owner.
Dessen Feedback (\json{} statt \textsc{csv}, englische Labels) wurde direkt in den
Sprint-2-Backlog überführt.

\paragraph{Reagieren auf Veränderung:}
Das Team akzeptierte die Änderungswünsche nach Sprint~1 ohne Mehrkosten-Diskussion und
priorisierte entsprechend im Sprint~2.

%% ============================================================
%%  KAPITEL 8: SCRUM-KONZEPT
%% ============================================================
\chapter{Scrum-Konzept}
\label{chap:scrum}

\section{Grundlagen}

Scrum ist ein leichtgewichtiges Rahmenwerk zur Entwicklung, Lieferung und Aufrechterhaltung
komplexer Produkte~\autocite{schwaber2020}. Es basiert auf drei Säulen:
\textbf{Transparenz}, \textbf{Inspektion} und \textbf{Adaption}.

Scrum definiert fünf Ereignisse, drei Artefakte und drei Rollen:

\begin{figure}[H]
  \centering
  \begin{tikzpicture}[
    box/.style={rectangle, rounded corners=3pt, minimum width=2.8cm, minimum height=0.8cm,
                text centered, font=\footnotesize\bfseries, draw=hsPrimary, thick},
    arr/.style={-{Stealth}, thick, color=hsPrimary},
    lbl/.style={font=\scriptsize, color=gray}
  ]
    \node[box, fill=hsPrimary!15] (pb)   at (0,0)     {Product Backlog};
    \node[box, fill=hsPrimary!30] (sp)   at (3.8,0)   {Sprint Planning};
    \node[box, fill=hsPrimary!45] (sb)   at (7.6,0)   {Sprint Backlog};
    \node[box, fill=hsPrimary!60, text=white] (sprint) at (7.6,-2) {Sprint (1--4 Wo.)};
    \node[box, fill=hsPrimary!45] (inc)  at (3.8,-2)  {Inkrement};
    \node[box, fill=hsPrimary!30] (rev)  at (0,-2)    {Review \& Retro};

    \draw[arr] (pb)  -- (sp);
    \draw[arr] (sp)  -- (sb);
    \draw[arr] (sb)  -- (sprint);
    \draw[arr] (sprint) -- (inc);
    \draw[arr] (inc) -- (rev);
    \draw[arr] (rev.north) |- ([yshift=0.6cm]pb.north) -| (pb.north);

    \node[lbl, above=0.1cm of pb]   {Priorisierung};
    \node[lbl, above=0.1cm of sp]   {Selektion};
    \node[lbl, above=0.1cm of sb]   {Aufgabenplanung};
    \node[lbl, right=0.1cm of sprint] {Daily Scrum};
    \node[lbl, below=0.1cm of inc]  {Nutzerwertt};
    \node[lbl, below=0.1cm of rev]  {Feedback};
  \end{tikzpicture}
  \caption{Scrum-Zyklus mit Artefakten und Ereignissen}
  \label{fig:scrum}
\end{figure}

\section{Rollen}

\begin{description}[leftmargin=2em, style=nextline]
  \item[Product Owner:] Verantwortlich für den maximalen Produktwert. Verwaltet und priorisiert
    das Product Backlog. Im Praktikum: Lehrveranstaltungsleitung.

  \item[Scrum Master:] Dienende Führungskraft, die Hindernisse beseitigt und den Scrum-Prozess
    schützt. Im Praktikum: Denise L.\ Doumbouya.

  \item[Development Team:] Selbstorganisiertes, cross-funktionales Team mit voller
    Implementierungsverantwortung. Im Praktikum: alle drei Gruppenmitglieder.
\end{description}

\section{Sprint-Länge und Ereignisse im Praktikum}

Aufgrund des Semesterrahmens wurden Sprints auf jeweils \textbf{eine Woche} festgelegt.
Daily Scrums fanden täglich als Kurzabstimmung (max.\ 15~Minuten) statt. Sprint Planning,
Sprint Review und Sprint Retrospektive folgten dem Scrum-Leitfaden, wurden jedoch auf
die praktischen Rahmenbedingungen der Lehrveranstaltung angepasst.

%% ============================================================
%%  KAPITEL 9: PRODUCT BACKLOG
%% ============================================================
\chapter{Product Backlog}
\label{chap:backlog}

\section{Initialer Product Backlog}

Der initiale Product Backlog wurde gemeinsam mit dem Product Owner in einer
Backlog-Refinement-Session erarbeitet. Die Priorisierung erfolgte nach dem
\textit{MoSCoW}-Verfahren: Must-have, Should-have, Could-have, Won't-have.

\begin{table}[H]
  \centering
  \caption{Product Backlog (nach erstem Refinement)}
  \label{tab:backlog}
  \begin{tabularx}{\textwidth}{C{0.7cm} L{5.5cm} C{2cm} C{2cm} C{1.5cm}}
    \toprule
    \textbf{ID} & \textbf{User Story} & \textbf{Prio.} & \textbf{Story Points} & \textbf{Sprint} \\
    \midrule
    US-01 & Als Lagermitarbeiter möchte ich neue Artikel anlegen, damit der Bestand vollständig ist.
           & Must & 3 & S1 \\
    US-02 & Als Lagermitarbeiter möchte ich alle Artikel anzeigen, damit ich den Überblick behalte.
           & Must & 2 & S1 \\
    US-03 & Als Lagermitarbeiter möchte ich den Bestand erhöhen, damit Wareneingänge erfasst werden.
           & Must & 2 & S1 \\
    US-04 & Als Lagermitarbeiter möchte ich den Bestand verringern, damit Warenabgänge erfasst werden.
           & Must & 2 & S1 \\
    US-05 & Als Manager möchte ich gewarnt werden, wenn Artikel den Mindestbestand unterschreiten.
           & Must & 3 & S1 \\
    US-06 & Als System-Administrator möchte ich den Bestand als Datei exportieren.
           & Should & 5 & S2 \\
    US-07 & Als internationaler Nutzer möchte ich englische Menütexte verwenden.
           & Should & 2 & S2 \\
    US-08 & Als Entwickler möchte ich eine saubere OOP-Struktur, damit der Code erweiterbar bleibt.
           & Should & 4 & S2 \\
    US-09 & Als Nutzer möchte ich fehlerhafte Eingaben abfangen, um Abstürze zu vermeiden.
           & Must & 3 & S1 \\
    \midrule
    & \textbf{Gesamt} & & \textbf{26} & \\
    \bottomrule
  \end{tabularx}
\end{table}

\section{Backlog-Refinement nach Sprint~1}

Nach dem Sprint Review von Sprint~1 wurde der Backlog substanziell überarbeitet.
Der Product Owner ergänzte drei neue Items und erhöhte die Priorität des JSON-Exports:

\begin{table}[H]
  \centering
  \caption{Backlog-Änderungen nach Sprint-1-Review}
  \label{tab:backlog-changes}
  \begin{tabularx}{\textwidth}{C{0.7cm} L{5.5cm} L{4.5cm} C{1.8cm}}
    \toprule
    \textbf{ID} & \textbf{Änderung} & \textbf{Begründung (Product Owner)} & \textbf{Aktion} \\
    \midrule
    US-06 & CSV \textrightarrow{} JSON als Export-Format & Standardformat; besser integrierbar & Angepasst \\
    US-07 & Priorität Must (war Should) & Internationaler Rollout geplant & Eskaliert \\
    US-08 & Klassen: Model/Service/UI-Trennung & Wartbarkeit für zukünftige Entwickler & Konkretisiert \\
    US-10 & Bestands-Summary mit Gesamtwert & Kaufmännische Kontrolle & Neu \\
    \bottomrule
  \end{tabularx}
\end{table}

%% ============================================================
%%  KAPITEL 10: SPRINTPLANUNG
%% ============================================================
\chapter{Sprintplanung}
\label{chap:sprintplanung}

\section{Sprint 1 \textendash{} Planung}

\textbf{Sprint-Ziel:} \textit{Ein lauffähiges Konsolenmenü mit vollständiger Artikelverwaltung
und Bestandsänderung ist fertiggestellt und testbar.}

Das Team nahm im Sprint Planning 16 Story Points aus dem Product Backlog in den Sprint~1-Backlog
auf. Die Velocity wurde für den ersten Sprint konservativ geschätzt, da das Team noch keine
gemeinsame Arbeitsgeschwindigkeit kannte.

\begin{table}[H]
  \centering
  \caption{Sprint-1-Backlog mit Aufgabenverteilung}
  \label{tab:sprint1}
  \begin{tabularx}{\textwidth}{C{0.7cm} L{4.5cm} L{3.5cm} C{1.5cm} C{2.5cm}}
    \toprule
    \textbf{ID} & \textbf{Aufgabe} & \textbf{Verantwortung} & \textbf{SP} & \textbf{Status (EOD)} \\
    \midrule
    T-01 & Klasse \texttt{Item} modellieren    & Jonas        & 2 & \cmark{} Done \\
    T-02 & Klasse \texttt{Warehouse} modellieren & Jonas      & 3 & \cmark{} Done \\
    T-03 & Konsolenmenü (\texttt{ConsoleUI})   & Denise       & 3 & \cmark{} Done \\
    T-04 & Bestand erhöhen / verringern        & Jonas        & 2 & \cmark{} Done \\
    T-05 & Negativbestand-Guard                & Priya        & 2 & \cmark{} Done \\
    T-06 & Mindestbestandswarnung              & Priya        & 2 & \cmark{} Done \\
    T-07 & Eingabe-Validierung (try-catch)     & Priya        & 2 & \cmark{} Done \\
    \midrule
    & \textbf{Gesamt}                          &              & \textbf{16} & \\
    \bottomrule
  \end{tabularx}
\end{table}

\subsection{Technische Entscheidungen in Sprint 1}

Im Sprint Planning wurde bewusst entschieden, die Klasse \texttt{Warehouse} als
\texttt{ArrayList<Item>} zu implementieren. Eine \texttt{HashMap}-basierte Struktur
wurde diskutiert, aber verworfen, da Iterationsreihenfolge für die Konsolenausgabe
wichtiger war als O(1)-Zugriff per Schlüssel.

Außerdem wurde festgelegt, die Menüsteuerung in einer separaten \texttt{ConsoleUI}-Klasse
zu kapseln, um die \texttt{Main}-Klasse schlank zu halten. Diese Entscheidung war in
Sprint~1 noch nicht konsequent umgesetzt, was im Review angemahnt wurde.

\section{Sprint 2 \textendash{} Planung}

\textbf{Sprint-Ziel:} \textit{Das System wird um JSON-Export, englische Labels und eine
klare Model-Service-UI-Architektur erweitert. Die Erkenntnisse aus dem Sprint-1-Review
sind vollständig umgesetzt.}

\begin{table}[H]
  \centering
  \caption{Sprint-2-Backlog mit Aufgabenverteilung}
  \label{tab:sprint2}
  \begin{tabularx}{\textwidth}{C{0.7cm} L{4.5cm} L{3.5cm} C{1.5cm} C{2.5cm}}
    \toprule
    \textbf{ID} & \textbf{Aufgabe} & \textbf{Verantwortung} & \textbf{SP} & \textbf{Status (EOD)} \\
    \midrule
    T-08 & \texttt{JsonExporter}-Klasse implementieren   & Priya  & 5 & \cmark{} Done \\
    T-09 & Gson-Dependency in \texttt{pom.xml} einbinden & Jonas  & 1 & \cmark{} Done \\
    T-10 & Alle Labels auf Englisch umstellen            & Denise & 2 & \cmark{} Done \\
    T-11 & \texttt{WarehouseService}-Klasse extrahieren  & Jonas  & 3 & \cmark{} Done \\
    T-12 & Refactoring \texttt{ConsoleUI}                & Denise & 2 & \cmark{} Done \\
    T-13 & Bestands-Summary (Gesamtwert)                 & Jonas  & 2 & \cmark{} Done \\
    T-14 & JUnit-Testfälle für Kernlogik                 & Priya  & 3 & \cmark{} Done \\
    \midrule
    & \textbf{Gesamt}                                    &        & \textbf{18} & \\
    \bottomrule
  \end{tabularx}
\end{table}

\subsection{Technische Entscheidungen in Sprint 2}

Die wohl wichtigste Architekturentscheidung in Sprint~2 war die Einführung der Klasse
\texttt{WarehouseService}. Im Wasserfall-Prototyp und noch in Sprint~1 war die
Geschäftslogik direkt in \texttt{Warehouse} (dem Datenmodell) verankert. Dieser Verstoß
gegen das \textit{Single Responsibility Principle} führte zu einem schwer testbaren,
eng gekoppelten Entwurf. Die Trennung in Datenklasse (\texttt{Item}), Repository
(\texttt{Warehouse}), Service (\texttt{WarehouseService}) und Präsentation
(\texttt{ConsoleUI}) folgte dem Prinzip der \textit{Separation of Concerns} und verbesserte
die Testbarkeit erheblich.

Für die JSON-Serialisierung wurde Googles \textsc{Gson}-Bibliothek gewählt, da sie keine
Annotationen erfordert, sehr leichtgewichtig ist und gut mit einfachen Java-POJOs
zusammenarbeitet.

%% ============================================================
%%  KAPITEL 11: SPRINT REVIEW
%% ============================================================
\chapter{Sprint Review}
\label{chap:review}

\section{Sprint-1-Review \textendash{} Protokoll}

\textbf{Datum:} Woche 2, Freitag, 14:00--14:45 Uhr\\
\textbf{Anwesende:} Gesamtes Entwicklungsteam, Product Owner (Lehrveranstaltungsleitung)\\
\textbf{Moderator:} Scrum Master Denise L.\ Doumbouya

\subsection{Präsentiertes Inkrement}

Das Team demonstrierte folgende funktionale Stories:

\begin{itemize}[leftmargin=1.5em]
  \item Start des Programmes mit vordefinierten Testartikeln (6 Artikel aus allen Kategorien)
  \item Anlegen eines neuen Artikels über das Konsolenmenü (US-01)
  \item Anzeige der Artikelliste (US-02) mit Bestands- und Preisangabe
  \item Bestand erhöhen (US-03) und verringern (US-04) mit Testfall für negativen Bestand
  \item Mindestbestandswarnung (US-05) für \texttt{Monitor} mit \texttt{quantity = 1, minimum = 3}
  \item Eingabe-Validierung bei ungültigem Menüpunkt (US-09)
\end{itemize}

\subsection{Kundenfeedback \textendash{} Product-Owner-Protokoll}

\begin{infobox}[Product Owner Feedback \textendash{} Sprint 1 Review]
  \textbf{Positiv:}
  \begin{itemize}[leftmargin=1.2em, itemsep=2pt]
    \item Kernfunktionalität vollständig und korrekt implementiert
    \item Mindestbestandswarnung funktioniert zuverlässig
    \item Eingabevalidierung verhindert Abstürze
  \end{itemize}

  \textbf{Verbesserungspotenzial:}
  \begin{itemize}[leftmargin=1.2em, itemsep=2pt]
    \item \textbf{Format:} CSV-Export nicht gewünscht \textrightarrow{} \textbf{JSON bevorzugt},
          da das System an eine Web-API angebunden werden soll
    \item \textbf{Sprache:} Deutsche Labels nicht akzeptabel für internationalen Rollout
          \textrightarrow{} \textbf{vollständig englische Benutzeroberfläche}
    \item \textbf{Architektur:} \texttt{Main.java} zu monolithisch (über 200 Zeilen);
          \textbf{klarere Klassen\-trennung} gewünscht
    \item \textbf{Zusatz:} Bestands-Summary mit Gesamtbestandswert wünschenswert
  \end{itemize}
\end{infobox}

\subsection{Reaktion des Teams}

Der Scrum Master dokumentierte das Feedback unmittelbar im Product Backlog. Im direkt
anschließenden Sprint-Planning-Meeting für Sprint~2 wurden alle vier Feedback-Punkte als
Sprint-2-Items aufgenommen. Die CSV-Exportklasse, die in Sprint~1 ansatzweise vorhanden war,
wurde im Backlog als \textit{obsolete} markiert und nicht weiterentwickelt.

\begin{warningbox}
  \textbf{Kritische Reflexion:} Der CSV-Ansatz stellte sich als klassisches Beispiel für
  \textit{Entwicklung am Kunden vorbei} heraus. Im Wasserfall wäre dies eine teure Änderung
  gewesen; im Scrum-Kontext kostete die Korrektur lediglich zwei Stunden Priya's Sprint-Zeit.
  Genau das ist der Mehrwert iterativer Entwicklung.
\end{warningbox}

\section{Sprint-2-Review \textendash{} Protokoll}

\textbf{Datum:} Woche 4, Freitag, 14:00--14:30 Uhr\\
\textbf{Anwesende:} Gesamtes Entwicklungsteam, Product Owner

\subsection{Präsentiertes Inkrement}

Das Team demonstrierte das vollständig überarbeitete System:

\begin{itemize}[leftmargin=1.5em]
  \item JSON-Export: Datei \texttt{warehouse.json} enthält alle 6 Testprodukte korrekt serialisiert
  \item Englisches Menü: alle Ausgaben konsistent in Englisch (\textit{Add item, Show items,
        Increase stock, Decrease stock, Export to JSON, Exit})
  \item \texttt{WarehouseService}-Klasse mit klarer Zuständigkeit (Logik-Schicht)
  \item \texttt{ConsoleUI} als reine Präsentationsschicht
  \item Bestands-Summary mit Gesamtbestandswert in Euro
  \item Erfolgreiche JUnit-Tests für alle Kernoperationen
\end{itemize}

\begin{successbox}
  \textbf{Product Owner Abnahme:} Alle Sprint-1-Feedback-Punkte wurden vollständig adressiert.
  Das System entspricht nun den fachlichen Anforderungen. Spezielle Hervorhebung durch den
  Product Owner: die saubere Klassenstruktur erleichtere eine spätere REST-API-Anbindung erheblich.
\end{successbox}

\section{Sprint Retrospektive}

In der Sprint-Retrospektive nach Sprint~2 reflektierte das Team den Gesamtprozess:

\begin{table}[H]
  \centering
  \caption{Ergebnisse der Sprint-Retrospektive}
  \label{tab:retro}
  \begin{tabularx}{\textwidth}{L{2cm} L{5cm} L{6cm}}
    \toprule
    \textbf{Kategorie} & \textbf{Thema} & \textbf{Maßnahme} \\
    \midrule
    \textcolor{darkGreen}{\textbf{Keep}}
      & Daily Standup per Nachrichtengruppe war effizient
      & Für zukünftige Projekte beibehalten \\
    \textcolor{darkGreen}{\textbf{Keep}}
      & Frühe Architekturentscheidungen dokumentieren
      & Architecture Decision Records (ADR) einführen \\
    \textcolor{orange!80!black}{\textbf{Improve}}
      & Sprint-1-Architektur war noch zu monolithisch
      & OOP-Prinzipien früher einplanen \\
    \textcolor{orange!80!black}{\textbf{Improve}}
      & Story-Point-Schätzung war ungenau (T-08 unterschätzt)
      & Planning Poker für nächstes Projekt nutzen \\
    \textcolor{darkRed}{\textbf{Stop}}
      & Technische Schulden aus Zeitdruck akzeptieren
      & Refactoring-Tasks explizit als Sprint-Items aufnehmen \\
    \bottomrule
  \end{tabularx}
\end{table}

%% ============================================================
%%  KAPITEL 12: SYSTEMARCHITEKTUR
%% ============================================================
\chapter{Systemarchitektur}
\label{chap:architektur}

\section{Paketstruktur}

Das finale System folgt einer Drei-Schichten-Architektur, die als Maven-Projekt strukturiert ist:

\begin{verbatim}
elager/
├── pom.xml
├── warehouse.json           (generiert)
└── src/
    ├── main/java/de/hs/elager/
    │   ├── model/
    │   │   └── Item.java
    │   ├── repository/
    │   │   └── Warehouse.java
    │   ├── service/
    │   │   └── WarehouseService.java
    │   ├── ui/
    │   │   └── ConsoleUI.java
    │   ├── export/
    │   │   └── JsonExporter.java
    │   └── Main.java
    └── test/java/de/hs/elager/
        ├── WarehouseServiceTest.java
        └── ItemTest.java
\end{verbatim}

\section{Architekturprinzipien}

Die Architektur folgt drei explizit gewählten Prinzipien:

\begin{description}[leftmargin=2em, style=nextline]
  \item[Separation of Concerns:] Datenmodell, Geschäftslogik, Persistenz und Präsentation
    sind in separaten Paketen und Klassen untergebracht.

  \item[Single Responsibility Principle (SRP):] Jede Klasse hat genau eine Aufgabe. Zum
    Beispiel ist \texttt{JsonExporter} ausschließlich für die Serialisierung zuständig;
    \texttt{WarehouseService} implementiert ausschließlich die Geschäftsregeln.

  \item[Dependency Inversion (vereinfacht):] \texttt{ConsoleUI} ruft ausschließlich
    \texttt{WarehouseService}-Methoden auf und kennt weder \texttt{Warehouse} noch
    \texttt{JsonExporter} direkt.
\end{description}

%% ============================================================
%%  KAPITEL 13: UML-KLASSENDIAGRAMM
%% ============================================================
\chapter{UML-Klassendiagramm}
\label{chap:uml}

\section{Beschreibung des Klassendiagramms}

Das folgende Klassendiagramm beschreibt die finale Systemarchitektur nach Sprint~2.

\begin{figure}[H]
  \centering
  \begin{tikzpicture}[
    class/.style={
      rectangle, draw=hsPrimary, thick,
      minimum width=4cm, text centered,
      font=\small
    },
    attr/.style={font=\ttfamily\footnotesize, text width=4cm, align=left},
    rel/.style={-{open triangle 60}, thick, color=hsPrimary},
    comp/.style={-{diamond}, thick, color=hsPrimary},
    uses/.style={dashed, -{Stealth}, thick, color=gray}
  ]
    % Item
    \node[class, fill=hsAccent] (item) at (0,0) {\textbf{Item}};
    \node[below=-0.02cm of item, rectangle, draw=hsPrimary, thick,
          minimum width=4cm, font=\ttfamily\footnotesize, align=left,
          fill=white, text width=3.8cm] (itemattr) {
      \textendash{} id: int\\
      \textendash{} name: String\\
      \textendash{} quantity: int\\
      \textendash{} minimum: int\\
      \textendash{} price: double
    };
    \node[below=-0.02cm of itemattr, rectangle, draw=hsPrimary, thick,
          minimum width=4cm, font=\ttfamily\footnotesize, align=left,
          fill=white, text width=3.8cm] (itemmeth) {
      + getters/setters()\\
      + isLowStock(): bool\\
      + toString(): String
    };

    % Warehouse
    \node[class, fill=hsAccent] (wh) at (6,0) {\textbf{Warehouse}};
    \node[below=-0.02cm of wh, rectangle, draw=hsPrimary, thick,
          minimum width=4.2cm, font=\ttfamily\footnotesize, align=left,
          fill=white, text width=4cm] (whattr) {
      \textendash{} items: List\textlangle{}Item\textrangle{}\\
      \textendash{} nextId: int
    };
    \node[below=-0.02cm of whattr, rectangle, draw=hsPrimary, thick,
          minimum width=4.2cm, font=\ttfamily\footnotesize, align=left,
          fill=white, text width=4cm] (whmeth) {
      + addItem(Item)\\
      + getItems(): List\\
      + findById(int): Item
    };

    % WarehouseService
    \node[class, fill=hsAccent] (srv) at (3,-5) {\textbf{WarehouseService}};
    \node[below=-0.02cm of srv, rectangle, draw=hsPrimary, thick,
          minimum width=4.8cm, font=\ttfamily\footnotesize, align=left,
          fill=white, text width=4.6cm] (srvattr) {
      \textendash{} warehouse: Warehouse
    };
    \node[below=-0.02cm of srvattr, rectangle, draw=hsPrimary, thick,
          minimum width=4.8cm, font=\ttfamily\footnotesize, align=left,
          fill=white, text width=4.6cm] (srvmeth) {
      + createItem(...): Item\\
      + increaseStock(id, qty)\\
      + decreaseStock(id, qty)\\
      + checkLowStock(): List\\
      + getTotalValue(): double
    };

    % JsonExporter
    \node[class, fill=hsAccent] (json) at (9,-5) {\textbf{JsonExporter}};
    \node[below=-0.02cm of json, rectangle, draw=hsPrimary, thick,
          minimum width=3.8cm, font=\ttfamily\footnotesize, align=left,
          fill=white, text width=3.6cm] (jsonattr) {
      \textendash{} gson: Gson
    };
    \node[below=-0.02cm of jsonattr, rectangle, draw=hsPrimary, thick,
          minimum width=3.8cm, font=\ttfamily\footnotesize, align=left,
          fill=white, text width=3.6cm] (jsonmeth) {
      + export(Warehouse, path)\\
      \textendash{} buildJson(...): String
    };

    % ConsoleUI
    \node[class, fill=hsAccent] (ui) at (3,-9.5) {\textbf{ConsoleUI}};
    \node[below=-0.02cm of ui, rectangle, draw=hsPrimary, thick,
          minimum width=4.2cm, font=\ttfamily\footnotesize, align=left,
          fill=white, text width=4cm] (uiattr) {
      \textendash{} service: WarehouseService\\
      \textendash{} exporter: JsonExporter\\
      \textendash{} scanner: Scanner
    };
    \node[below=-0.02cm of uiattr, rectangle, draw=hsPrimary, thick,
          minimum width=4.2cm, font=\ttfamily\footnotesize, align=left,
          fill=white, text width=4cm] (uimeth) {
      + start()\\
      \textendash{} showMenu()\\
      \textendash{} handleInput(int)
    };

    % Beziehungen
    \draw[comp] (wh.west) -- node[above, font=\scriptsize]{1..*} (item.east);
    \draw[uses] (srv.north) -- (wh.south);
    \draw[uses] (json.north) -- (wh.south);
    \draw[uses] (ui.north) -- (srv.south);
    \draw[uses] ([xshift=1cm]ui.north) -- (json.south);
  \end{tikzpicture}
  \caption{UML-Klassendiagramm des Elektroniklagerverwaltungssystems (finaler Zustand)}
  \label{fig:uml}
\end{figure}

\section{Beschreibung der Klassen}

\begin{description}[leftmargin=2em, style=nextline]
  \item[\texttt{Item} (Modellschicht):] Reine Datenhaltungsklasse (\textit{Plain Old Java Object}).
    Enthält alle Artikelattribute sowie die Hilfsmethode \texttt{isLowStock()}, die prüft, ob
    \texttt{quantity < minimum} gilt.

  \item[\texttt{Warehouse} (Repository-Schicht):] Kapselt die Artikelliste und den
    ID-Zähler. Verantwortlich für das Hinzufügen neuer Artikel und die Suche per ID.
    Enthält keine Geschäftslogik.

  \item[\texttt{WarehouseService} (Service-Schicht):] Implementiert alle Geschäftsregeln:
    Negativbestand-Guard, Mindestbestandsprüfung, Gesamtwertberechnung. Enthält keine
    Ausgabe-Logik.

  \item[\texttt{JsonExporter} (Persistenz-Schicht):] Kapselt die Gson-Bibliothek und
    serialisiert die Artikelliste in eine \json{}-Datei. Fehler beim Schreiben werden als
    \texttt{RuntimeException} propagiert.

  \item[\texttt{ConsoleUI} (Präsentationsschicht):] Liest Benutzereingaben, delegiert
    an \texttt{WarehouseService} und \texttt{JsonExporter}, gibt Ergebnisse auf der Konsole aus.

  \item[\texttt{Main}:] Einstiegspunkt. Instanziiert alle Objekte und startet
    \texttt{ConsoleUI.start()}.
\end{description}

%% ============================================================
%%  KAPITEL 14: IMPLEMENTIERUNG
%% ============================================================
\chapter{Implementierung}
\label{chap:implementierung}

\section{Klasse \texttt{Item}}

\begin{lstlisting}[caption={Item.java -- Datenmodell eines Lagerartikels}, label={lst:item}]
package de.hs.elager.model;

public class Item {

    private int    id;
    private String name;
    private int    quantity;
    private int    minimum;
    private double price;

    public Item(int id, String name, int quantity, int minimum, double price) {
        this.id       = id;
        this.name     = name;
        this.quantity = quantity;
        this.minimum  = minimum;
        this.price    = price;
    }

    public boolean isLowStock() {
        return this.quantity < this.minimum;
    }

    // --- Getters & Setters ---
    public int    getId()       { return id; }
    public String getName()     { return name; }
    public int    getQuantity() { return quantity; }
    public int    getMinimum()  { return minimum; }
    public double getPrice()    { return price; }

    public void setQuantity(int quantity) { this.quantity = quantity; }

    @Override
    public String toString() {
        return String.format("[%d] %-25s | Qty: %3d | Min: %2d | Price: %8.2f EUR",
                id, name, quantity, minimum, price);
    }
}
\end{lstlisting}

\section{Klasse \texttt{WarehouseService}}

\begin{lstlisting}[caption={WarehouseService.java -- Geschäftslogik}, label={lst:service}]
package de.hs.elager.service;

import de.hs.elager.model.Item;
import de.hs.elager.repository.Warehouse;
import java.util.List;
import java.util.stream.Collectors;

public class WarehouseService {

    private final Warehouse warehouse;

    public WarehouseService(Warehouse warehouse) {
        this.warehouse = warehouse;
    }

    public Item createItem(String name, int quantity, int minimum, double price) {
        Item item = new Item(warehouse.getNextId(), name, quantity, minimum, price);
        warehouse.addItem(item);
        return item;
    }

    public void increaseStock(int id, int amount) {
        if (amount <= 0) throw new IllegalArgumentException("Amount must be positive.");
        Item item = warehouse.findById(id);
        item.setQuantity(item.getQuantity() + amount);
    }

    public void decreaseStock(int id, int amount) {
        if (amount <= 0) throw new IllegalArgumentException("Amount must be positive.");
        Item item = warehouse.findById(id);
        int newQty = item.getQuantity() - amount;
        if (newQty < 0) {
            throw new IllegalStateException(
                "Cannot decrease stock below zero. Current: "
                + item.getQuantity() + ", Requested: " + amount);
        }
        item.setQuantity(newQty);
    }

    public List<Item> getLowStockItems() {
        return warehouse.getItems().stream()
                .filter(Item::isLowStock)
                .collect(Collectors.toList());
    }

    public double getTotalValue() {
        return warehouse.getItems().stream()
                .mapToDouble(i -> i.getQuantity() * i.getPrice())
                .sum();
    }

    public List<Item> getAllItems() {
        return warehouse.getItems();
    }
}
\end{lstlisting}

\section{Klasse \texttt{JsonExporter}}

\begin{lstlisting}[caption={JsonExporter.java -- JSON-Serialisierung}, label={lst:json}]
package de.hs.elager.export;

import com.google.gson.Gson;
import com.google.gson.GsonBuilder;
import de.hs.elager.repository.Warehouse;
import java.io.FileWriter;
import java.io.IOException;

public class JsonExporter {

    private final Gson gson;

    public JsonExporter() {
        this.gson = new GsonBuilder().setPrettyPrinting().create();
    }

    public void export(Warehouse warehouse, String filePath) {
        try (FileWriter writer = new FileWriter(filePath)) {
            gson.toJson(warehouse.getItems(), writer);
            System.out.println("Export successful: " + filePath);
        } catch (IOException e) {
            throw new RuntimeException("JSON export failed: " + e.getMessage(), e);
        }
    }
}
\end{lstlisting}

\section{Klasse \texttt{ConsoleUI}}

\begin{lstlisting}[caption={ConsoleUI.java -- Präsentationsschicht (Auszug)}, label={lst:ui}]
package de.hs.elager.ui;

import de.hs.elager.export.JsonExporter;
import de.hs.elager.model.Item;
import de.hs.elager.service.WarehouseService;
import java.util.List;
import java.util.Scanner;

public class ConsoleUI {

    private final WarehouseService service;
    private final JsonExporter     exporter;
    private final Scanner          scanner;

    public ConsoleUI(WarehouseService service, JsonExporter exporter) {
        this.service  = service;
        this.exporter = exporter;
        this.scanner  = new Scanner(System.in);
    }

    public void start() {
        boolean running = true;
        while (running) {
            showMenu();
            int choice = readInt();
            running = handleInput(choice);
        }
        System.out.println("Goodbye.");
    }

    private void showMenu() {
        System.out.println("\n=== ELECTRONICS WAREHOUSE MANAGEMENT ===");
        System.out.println("1. Add item");
        System.out.println("2. Show all items");
        System.out.println("3. Increase stock");
        System.out.println("4. Decrease stock");
        System.out.println("5. Export to JSON");
        System.out.println("6. Show inventory summary");
        System.out.println("0. Exit");
        System.out.print("Choose: ");
    }

    private boolean handleInput(int choice) {
        switch (choice) {
            case 1  -> addItem();
            case 2  -> showItems();
            case 3  -> changeStock(true);
            case 4  -> changeStock(false);
            case 5  -> exportJson();
            case 6  -> showSummary();
            case 0  -> { return false; }
            default -> System.out.println("Invalid option. Please try again.");
        }
        return true;
    }

    private void showItems() {
        List<Item> items = service.getAllItems();
        if (items.isEmpty()) { System.out.println("No items in warehouse."); return; }
        System.out.println("\n--- ITEMS ---");
        items.forEach(System.out::println);
        List<Item> lowStock = service.getLowStockItems();
        if (!lowStock.isEmpty()) {
            System.out.println("\n[WARNING] Low stock items:");
            lowStock.forEach(i -> System.out.println("  >> " + i.getName()
                + " (qty=" + i.getQuantity() + ", min=" + i.getMinimum() + ")"));
        }
    }

    private void changeStock(boolean increase) {
        System.out.print("Item ID: ");
        int id = readInt();
        System.out.print("Amount: ");
        int amount = readInt();
        try {
            if (increase) service.increaseStock(id, amount);
            else          service.decreaseStock(id, amount);
            System.out.println("Stock updated successfully.");
        } catch (IllegalStateException | IllegalArgumentException e) {
            System.out.println("[ERROR] " + e.getMessage());
        }
    }

    private int readInt() {
        while (!scanner.hasNextInt()) {
            System.out.print("Invalid input. Enter a number: ");
            scanner.next();
        }
        return scanner.nextInt();
    }
    // ... weitere Methoden (addItem, exportJson, showSummary)
}
\end{lstlisting}

\section{Beispiel-Output: JSON-Export}

Das folgende Listing zeigt einen Ausschnitt der generierten Datei \texttt{warehouse.json}:

\begin{lstlisting}[language={}, caption={warehouse.json -- Beispielausgabe (Auszug)}, label={lst:jsonout}]
[
  {
    "id": 1,
    "name": "Dell UltraSharp 27\" 4K Monitor",
    "quantity": 5,
    "minimum": 3,
    "price": 599.0
  },
  {
    "id": 2,
    "name": "Samsung 980 Pro 1TB SSD",
    "quantity": 12,
    "minimum": 5,
    "price": 89.99
  },
  {
    "id": 3,
    "name": "ASUS RT-AX88U WiFi 6 Router",
    "quantity": 2,
    "minimum": 4,
    "price": 229.0
  }
]
\end{lstlisting}

%% ============================================================
%%  KAPITEL 15: DEFINITION OF DONE
%% ============================================================
\chapter{Definition of Done}
\label{chap:dod}

Die \dod{} (\textsc{DoD}) ist ein explizites Qualitätskriterium, das festlegt, wann eine
User Story als vollständig fertiggestellt gilt. Sie wurde im Sprint Planning von Sprint~1
gemeinsam mit dem Product Owner vereinbart und in Sprint~2 um Testanforderungen erweitert.

\begin{table}[H]
  \centering
  \caption{Definition of Done \textendash{} Gültig ab Sprint 2}
  \label{tab:dod}
  \begin{tabularx}{\textwidth}{C{0.5cm} L{8cm} C{1.5cm} C{1.5cm}}
    \toprule
    \textbf{\#} & \textbf{Kriterium} & \textbf{S1} & \textbf{S2} \\
    \midrule
    1 & Code ist kompilierbar und ausführbar & \cmark & \cmark \\
    2 & Alle akzeptierten Akzeptanzkriterien der User Story sind erfüllt & \cmark & \cmark \\
    3 & Mindestens ein positiver und ein negativer Testfall existiert & \xmark & \cmark \\
    4 & Keine offenen \texttt{NullPointerException}-Risiken & \cmark & \cmark \\
    5 & Kein Code verletzt das Single Responsibility Principle & \xmark & \cmark \\
    6 & Alle Variablen- und Methodennamen sind englisch & \xmark & \cmark \\
    7 & Code ist vom zweiten Teammitglied reviewed (4-Augen-Prinzip) & \cmark & \cmark \\
    8 & Änderungen sind ins Git-Repository committet (aussagekräftige Message) & \cmark & \cmark \\
    9 & Fehlerbehandlung für alle Benutzereingaben implementiert & \cmark & \cmark \\
    10 & JSON-Export erzeugt valide Datei (Sprint 2 ab US-06) & n.a. & \cmark \\
    \bottomrule
  \end{tabularx}
\end{table}

\begin{infobox}[Bedeutung der DoD für das Praktikum]
  Die DoD verhinderte in Sprint~2 eine vorzeitige Fertigstellung von US-07 (englische Labels):
  Jonas hatte die Menüpunkte übersetzt, aber die Fehlermeldungen innerhalb der Geschäftslogik
  waren noch auf Deutsch. Das 4-Augen-Prinzip im Code Review identifizierte dieses Problem,
  bevor die Story als \textit{done} akzeptiert wurde.
\end{infobox}

%% ============================================================
%%  KAPITEL 16: TESTKONZEPT
%% ============================================================
\chapter{Testkonzept}
\label{chap:tests}

\section{Teststrategie}

Das Testkonzept unterscheidet zwischen manuellen Funktionstests (Sprint~1) und automatisierten
JUnit-Tests (Sprint~2). Diese Progression spiegelt die iterative Qualitätssteigerung wider.

\section{Manuelle Testfälle (Sprint 1)}

\begin{longtable}{C{0.7cm} L{4cm} L{3.5cm} L{3.5cm} C{1.5cm}}
  \caption{Manuelle Testfälle aus Sprint 1}
  \label{tab:tests-manual} \\
  \toprule
  \textbf{TC} & \textbf{Beschreibung} & \textbf{Eingabe} & \textbf{Erwartet} & \textbf{Ergebnis} \\
  \midrule
  \endfirsthead
  \multicolumn{5}{c}{\tablename~\thetable{} (Fortsetzung)} \\
  \toprule
  \textbf{TC} & \textbf{Beschreibung} & \textbf{Eingabe} & \textbf{Erwartet} & \textbf{Ergebnis} \\
  \midrule
  \endhead
  \bottomrule
  \endlastfoot
  MT-01 & Artikel anlegen & Name: ``Keyboard'', qty: 10, min: 3, price: 49.99 & Artikel in Liste sichtbar mit ID~1 & \cmark \\
  MT-02 & Bestand erhöhen & ID~1, Amount: 5 & quantity = 15 & \cmark \\
  MT-03 & Bestand verringern (gültig) & ID~1, Amount: 3 & quantity = 12 & \cmark \\
  MT-04 & Bestand verringern (negativ) & ID~1, Amount: 99 & Fehlermeldung, qty unverändert & \cmark \\
  MT-05 & Mindestbestandswarnung & qty = 2, min = 5 & Warnung ausgegeben & \cmark \\
  MT-06 & Ungültige Menüeingabe & Text statt Zahl & Retry-Aufforderung & \cmark \\
  MT-07 & Ungültige ID beim Bestand ändern & ID~999 & Fehlermeldung & \cmark \\
\end{longtable}

\section{JUnit-Testfälle (Sprint 2)}

\begin{lstlisting}[caption={WarehouseServiceTest.java -- JUnit-5-Tests}, label={lst:tests}]
package de.hs.elager.service;

import de.hs.elager.model.Item;
import de.hs.elager.repository.Warehouse;
import org.junit.jupiter.api.*;
import static org.junit.jupiter.api.Assertions.*;

class WarehouseServiceTest {

    private WarehouseService service;

    @BeforeEach
    void setUp() {
        Warehouse wh = new Warehouse();
        service = new WarehouseService(wh);
        service.createItem("Test Monitor", 10, 3, 299.99);
    }

    @Test
    @DisplayName("createItem increases warehouse size by 1")
    void testCreateItem() {
        int before = service.getAllItems().size();
        service.createItem("Test SSD", 5, 2, 79.99);
        assertEquals(before + 1, service.getAllItems().size());
    }

    @Test
    @DisplayName("increaseStock adds correct amount")
    void testIncreaseStock() {
        service.increaseStock(1, 5);
        assertEquals(15, service.getAllItems().get(0).getQuantity());
    }

    @Test
    @DisplayName("decreaseStock subtracts correct amount")
    void testDecreaseStock() {
        service.decreaseStock(1, 4);
        assertEquals(6, service.getAllItems().get(0).getQuantity());
    }

    @Test
    @DisplayName("decreaseStock throws on negative result")
    void testDecreaseStockNegative() {
        assertThrows(IllegalStateException.class,
            () -> service.decreaseStock(1, 99));
    }

    @Test
    @DisplayName("getLowStockItems returns item when below minimum")
    void testLowStock() {
        service.decreaseStock(1, 8); // qty=2, min=3 -> low stock
        assertFalse(service.getLowStockItems().isEmpty());
        assertEquals("Test Monitor", service.getLowStockItems().get(0).getName());
    }

    @Test
    @DisplayName("getTotalValue returns correct sum")
    void testTotalValue() {
        double expected = 10 * 299.99;
        assertEquals(expected, service.getTotalValue(), 0.001);
    }

    @Test
    @DisplayName("increaseStock throws on non-positive amount")
    void testIncreaseStockInvalidAmount() {
        assertThrows(IllegalArgumentException.class,
            () -> service.increaseStock(1, -3));
    }
}
\end{lstlisting}

\section{Testergebnisse Sprint 2}

\begin{table}[H]
  \centering
  \caption{JUnit-Testergebnisse nach Sprint~2}
  \label{tab:testergebnisse}
  \begin{tabularx}{0.8\textwidth}{L{6cm} C{2cm} C{2cm}}
    \toprule
    \textbf{Testklasse} & \textbf{Tests} & \textbf{Ergebnis} \\
    \midrule
    \texttt{WarehouseServiceTest}  & 7 & \textcolor{darkGreen}{\textbf{7/7 PASS}} \\
    \texttt{ItemTest}              & 3 & \textcolor{darkGreen}{\textbf{3/3 PASS}} \\
    \midrule
    \textbf{Gesamt}                & \textbf{10} & \textcolor{darkGreen}{\textbf{10/10 PASS}} \\
    \bottomrule
  \end{tabularx}
\end{table}

%% ============================================================
%%  KAPITEL 17: VERGLEICH WASSERFALL VS. SCRUM
%% ============================================================
\chapter{Vergleich: Wasserfallmodell vs.\ Scrum}
\label{chap:vergleich}

\section{Struktureller Vergleich}

\begin{table}[H]
  \centering
  \caption{Direktvergleich: Wasserfallmodell vs.\ Scrum}
  \label{tab:vergleich}
  \begin{tabularx}{\textwidth}{L{3.5cm} L{5cm} L{5cm}}
    \toprule
    \textbf{Kriterium} & \textbf{Wasserfallmodell} & \textbf{Scrum} \\
    \midrule
    Planungsphilosophie & Vollständige Planung vor Start & Just-Enough-Planning, iterativ \\
    Flexibilität & Gering; Änderungen kosten viel & Hoch; Änderungen nach jedem Sprint \\
    Kundeninteraktion & Anfang und Ende & Kontinuierlich (Sprint Review) \\
    Risiko-Timing & Risiken zeigen sich spät & Risiken früh sichtbar \\
    Lieferzeitpunkt & Einmaliger Release am Ende & Inkrement nach jedem Sprint \\
    Dokumentation & Umfangreich, vorab & Minimal, begleitend \\
    Teamstruktur & Hierarchisch / funktional & Selbstorganisiert, cross-funktional \\
    Geeignet für & Stabile Anforderungen, hohe Compliance & Veränderliche Anforderungen \\
    Messbarkeit & Meilensteine & Velocity, Burndown-Chart \\
    Fehlerkosten & Steigen exponentiell mit Phasenverzug & Begrenzt auf Sprint-Umfang \\
    \bottomrule
  \end{tabularx}
\end{table}

\section{Projektspezifischer Vergleich}

Im Rahmen dieses Praktikums manifestierten sich die theoretischen Unterschiede sehr konkret:

\begin{table}[H]
  \centering
  \caption{Vergleich am konkreten Projektbeispiel}
  \label{tab:projektvergleich}
  \begin{tabularx}{\textwidth}{L{3.5cm} L{5cm} L{5cm}}
    \toprule
    \textbf{Situation} & \textbf{Wasserfall} & \textbf{Scrum} \\
    \midrule
    Exportformat-Entscheidung &
      CSV fest im Pflichtenheft verankert; Änderung erfordert Neuplanung &
      JSON nach Sprint~1-Review ohne Prozessbruch in Sprint~2 \\
    Englische Labels &
      Deutsch im Pflichtenheft; keine Feedback-Schleife &
      Nach Review sofort priorisiert und umgesetzt \\
    Architektur &
      Monolithische \texttt{Main}-Klasse im Design festgelegt &
      Refactoring als Sprint-Item nach Review \\
    Qualitätssicherung &
      Nur manuelle Tests am Projektende &
      JUnit-Tests ab Sprint~2 als DoD-Kriterium \\
    \bottomrule
  \end{tabularx}
\end{table}

\section{Eignung nach Projekttyp}

\begin{description}[leftmargin=2em, style=nextline]
  \item[Wasserfall eignet sich für:]
    Projekte mit streng reguliertem Umfeld (Medizintechnik, Luft- und Raumfahrt,
    Embedded-Systeme mit Zertifizierungspflicht), sehr stabilen und vollständig
    bekannten Anforderungen, festem Budget und fixiertem Lieferumfang sowie
    Projekten, bei denen Änderungen physisch teuer sind (Hardwareentwicklung).

  \item[Scrum eignet sich für:]
    Projekte mit sich entwickelnden oder noch nicht vollständig bekannten Anforderungen,
    enger Kundeninteraktion während der Entwicklung, Digitalisierungs- und
    Produktentwicklungsvorhaben, Start-up-Kontexten und Inhouse-Entwicklung sowie
    alle Situationen, in denen frühe Nutzerfeedbacks den Produktwert erheblich steigern können.
\end{description}

%% ============================================================
%%  KAPITEL 18: REFLEXION
%% ============================================================
\chapter{Reflexion}
\label{chap:reflexion}

\section{Einleitendes}

Die folgende Reflexion ist der Kern dieser Dokumentation. Sie geht über eine
Zusammenfassung des Geleisteten hinaus und analysiert ehrlich, an welchen Stellen die
Theorie auf die Realität des Praktikums traf \textendash{} und an welchen sie scheiterte.

\section{In welchem Modell fühlte sich das Team sicherer?}

Eine bemerkenswert ehrliche Antwort auf diese Frage: Im Wasserfallmodell fühlte sich das
Team anfänglich \textit{sicherer} \textendash{} aber nur im Sinne einer trügerischen Sicherheit.

Die klaren Phasenübergänge, das abgeschlossene Pflichtenheft und die vorgegebene Sequenz
erzeugten das Gefühl, \textit{alles unter Kontrolle} zu haben. Jeder wusste genau, was zu tun
war, wann es getan sein musste und in welcher Reihenfolge. Diese Klarheit ist psychologisch
angenehm, insbesondere für ein Team ohne gemeinsame Vorerfahrung.

Im Scrum-Kontext war diese Planungssicherheit zunächst abwesend. Das Sprint Planning
mit relativen Story-Point-Schätzungen wirkte anfangs willkürlich. Die Frage
\textit{Wann sind wir fertig?} ließ sich nicht mehr mit einem Datum beantworten, sondern
nur noch mit \textit{Am Ende von Sprint 2 \textendash{} wenn alles Geplante die DoD erfüllt}.

Mit dem Ende des zweiten Sprints hatte sich die Wahrnehmung grundlegend verschoben:
Die Sicherheit des Wasserfall-Ansatzes entpuppte sich als Illusion. Der Wasserfall-Prototyp
lieferte zwar planmäßig, aber am falschen Ziel vorbei. Das Scrum-Inkrement nach Sprint~2
hingegen war trotz anfänglicher Orientierungslosigkeit \textbf{exakt das, was der Product
Owner gebraucht hatte}.

\section{Wann wurde deutlich, dass am Kunden vorbei entwickelt wurde?}

Der Moment war präzise datierbar: der Sprint-1-Review, Freitag der zweiten
Praktikumswoche. Als der Product Owner das laufende System sah, fiel der erste Kommentar
nicht auf die Funktionalität \textendash{} die war korrekt \textendash{} sondern auf das Format: \textit{Warum
CSV? Wir brauchen JSON.}

Dieser Satz war lehrreich auf mehreren Ebenen. Erstens bestätigte er, dass
Anforderungen, die auf dem Papier eindeutig erscheinen (\textit{``Exportfunktion''
im Pflichtenheft), in der Vorstellung des Auftraggebers sehr konkret sind \textendash{} und
diese Konkretheit kommt oft erst beim Anblick eines lauffähigen Systems zutage.
Zweitens offenbarte er ein strukturelles Problem des Wasserfall-Modells: Das Pflichtenheft
hatte keine Zeile über das Exportformat gesagt, weil in der Anforderungsphase niemand
danach gefragt hatte. Im Wasserfall hätte das eine Vertragsänderung bedeutet.

Die zweite Erkenntnis war subtiler: Die deutschen Labels. Keines der Teammitglieder
hatte beim Schreiben des deutschen Menüs daran gedacht, den Product Owner zu fragen, ob
die Sprache relevant sei. Im Pflichtenheft stand \textit{Konsolenmenü} \textendash{} nicht
\textit{englisches Konsolenmenü}. Das ist der Unterschied zwischen einem spezifizierten
System und einem \textit{nützlichen} System.

Im Wasserfall-Prototyp fehlte schlicht die Gelegenheit, diese Lücken zu schließen.
Im Scrum-Prozess wurde diese Erkenntnis zu einem produktiven Sprint-2-Item.

\section{Für welche Projekte eignet sich Wasserfall?}

Die reflexive Antwort nach diesem Praktikum lautet: Wasserfall eignet sich für Projekte,
bei denen das Risiko einer Anforderungsänderung bewusst und akzeptiert ausgeschlossen wird.

Dies ist keine Kritik an dem Modell, sondern eine Klarstellung seines Geltungsbereichs.
In sicherheitskritischen Systemen \textendash{} medizinische Software, avionische Steuerungen,
Industriesteuerungen nach IEC~61508 \textendash{} sind Änderungsprozesse bewusst formalisiert
und teuer, weil jede Änderung eine vollständige Requalifizierung nach sich ziehen kann.
Dort ist der Wasserfall nicht nur akzeptabel, sondern \textit{geboten}.

Für Projekte mit fest vereinbartem Lieferumfang und klar definierten
Qualitätsmerkmalen \textendash{} etwa Wartungsverträge, Migrationsprojekte mit stabiler
Datenbasis oder Batch-Verarbeitungssysteme für interne Prozesse \textendash{} ist das
Wasserfallmodell effizienter als Scrum, weil der Overhead der Scrum-Zeremonien nicht
durch erhöhte Flexibilität aufgewogen wird.

\section{Für welche Projekte eignet sich Scrum?}

Scrum entfaltet seinen vollen Wert dort, wo Anforderungen sich im Laufe der Entwicklung
konkretisieren oder verändern. Dieses Praktikum ist ein Lehrbuchbeispiel: Die Anforderung
\textit{Exportfunktion} klingt eindeutig, ist es aber nicht \textendash{} bis man den Kontext kennt.
Eben diesen Kontext liefert das Sprint Review.

Darüber hinaus eignet sich Scrum besonders dann, wenn das Entwicklungsteam nah am
zukünftigen Nutzer arbeitet und kontinuierliches Feedback integrieren kann. In der
kommerziellen Softwareentwicklung \textendash{} ERP-Systeme, CRM-Lösungen,
E-Commerce-Plattformen \textendash{} ist Scrum heute de facto der Standard, weil
Fachbereiche selten in der Lage sind, Anforderungen zu Beginn vollständig zu spezifizieren.

\section{Welche Rolle spielte das Sprint Review?}

Das Sprint Review war das wichtigste Ereignis des gesamten Praktikums. Nicht das
Sprint Planning, nicht die Daily Scrums, sondern das Review \textendash{} weil es der einzige
Moment war, in dem der Product Owner das System \textit{tatsächlich genutzt} hat.

Eine schriftliche Anforderung wird interpretiert. Ein laufendes System wird erlebt.
Dieser Unterschied ist fundamental. Im Sprint-1-Review erlebte der Product Owner, dass
CSV als Exportformat für seinen geplanten Verwendungszweck unbrauchbar war. Diese
Erkenntnis hätte keine noch so gründliche Anforderungsanalyse zu Projektbeginn zuverlässig
hervorgebracht, weil die Frage \textit{In welchem Format soll der Export erfolgen?}
in einem abstrakten Pflichtenheft-Gespräch anders beantwortet wird als beim Anblick
eines konkreten CSV-Exports.

Das Review erzeugte außerdem Motivation im Team. Positives Feedback des Product Owners
auf die funktionierende Kernfunktionalität war weitaus wirkungsvoller als jedes
formale Status-Meeting.

\section{Warum war Scrum bei diesem Projekt sinnvoll?}

Scrum war bei diesem Projekt sinnvoll, weil das System für einen realen (simulierten)
Auftraggeber entwickelt wurde, dessen vollständige Vorstellungen zu Beginn nicht bekannt
waren. Das ist keine Ausnahme \textendash{} es ist der Regelfall in der Softwareentwicklung.

Konkret profitierte das Projekt an drei Stellen von Scrum:

\begin{enumerate}[leftmargin=1.5em]
  \item \textbf{Feedback-Schleife:} Das Sprint Review erzeugte konkrete, umsetzbare
    Verbesserungspunkte. Ohne Review wäre Sprint~2 ohne Korrekturen geblieben.

  \item \textbf{Priorisierung:} Das Backlog-Refinement nach Sprint~1 erlaubte es,
    JSON-Export \textendash{} nun als Must-have klassifiziert \textendash{} sofort in Sprint~2
    zu priorisieren, ohne den Projektzeitplan zu gefährden.

  \item \textbf{Kontinuierliche Qualitätssteigerung:} Die DoD-Erweiterung zwischen
    Sprint~1 und Sprint~2 (Einführung der JUnit-Tests) führte zu messbarer
    Qualitätsverbesserung, die im Wasserfallmodell nicht möglich gewesen wäre, weil
    der Testzeitpunkt am Ende fixiert war.
\end{enumerate}

\section{Unterschiede zwischen Theorie und realem Praktikum}

Die akademische Scrum-Theorie beschreibt ein idealisiertes Framework. Die Realität des
Praktikums war an mehreren Stellen anders:

\begin{description}[leftmargin=2em, style=nextline]
  \item[Daily Scrum vs.\ Realität:] Tägliche 15-minütige Standups sind in einem
    dreistündigen Wochenpraktikum nicht sinnvoll umsetzbar. Stattdessen nutzte das Team
    eine Nachrichtengruppe für tägliche Statusmeldungen. Die Substanz des Daily Scrum
    \textendash{} Hindernisse sichtbar machen \textendash{} blieb erhalten; die Form wurde pragmatisch
    angepasst.

  \item[Velocity vs.\ Realität:] Story-Point-Schätzungen waren im ersten Sprint
    erheblich ungenauer als erwartet. Aufgabe T-08 (JSON-Export) wurde mit 5 Story Points
    geschätzt, benötigte aber in der Umsetzung deutlich mehr Zeit, da die Gson-Bibliothek
    zunächst korrekt in Maven eingebunden werden musste und die Pretty-Printing-Konfiguration
    mehr Einarbeitungszeit erforderte als geplant.

  \item[Product Owner vs.\ Realität:] Ein wirklicher Product Owner ist täglich verfügbar.
    Im Praktikum stand die Lehrveranstaltungsleitung nur im Sprint Review zur Verfügung.
    Dies führte dazu, dass Zwischenfragen nicht sofort beantwortet werden konnten, was
    in Sprint~1 zu einer kleinen Fehlentscheidung (Menü-Strukturierung) führte, die erst
    im Review korrigiert wurde.

  \item[Selbstorganisation vs.\ Realität:] Selbstorganisierte Teams benötigen ein
    Mindestmaß an gemeinsamer Erfahrung. Im ersten Sprint übernahm der Scrum Master
    \textendash{} entgegen der Scrum-Theorie \textendash{} implizit auch Koordinierungsaufgaben,
    die eigentlich dem Team gehören. Ab Sprint~2 hatte das Team genug gemeinsamen Rhythmus
    entwickelt, um echte Selbstorganisation zu praktizieren.
\end{description}

Diese Abweichungen schmälern den Erkenntniswert nicht \textendash{} im Gegenteil: Sie belegen,
dass Scrum kein starres Regelwerk ist, sondern ein Rahmenwerk, das bewusste Anpassung
an den Kontext erfordert.

%% ============================================================
%%  KAPITEL 19: AUTHENTISCHER PROJEKTABLAUF
%% ============================================================
\chapter{Authentischer Projektablauf}
\label{chap:ablauf}

\section{Chronologische Projekthistorie}

\begin{table}[H]
  \centering
  \caption{Projektzeitstrahl mit Meilensteinen}
  \label{tab:zeitstrahl}
  \begin{tabularx}{\textwidth}{C{1.5cm} L{3cm} L{8cm}}
    \toprule
    \textbf{Woche} & \textbf{Phase} & \textbf{Ereignisse \& Erkenntnisse} \\
    \midrule
    W1 & Wasserfall & Anforderungsanalyse, Pflichtenheft, Design (CSV, Deutsch) \\
    W1 & Wasserfall & Implementierung Grundgerüst (\texttt{Item}, \texttt{Warehouse}, \texttt{Main}) \\
    W2 & Wasserfall & Fertigstellung, manuelle Tests, Präsentation beim Product Owner \\
    W2 & Übergang & \textbf{Pivot:} Wechsel zu Scrum. Backlog erstellen, Rollen verteilen \\
    W2 & Sprint 1 & Planning: 16 SP, Sprint-Ziel vereinbart \\
    W2--W3 & Sprint 1 & Entwicklung: Kernfunktionen, Konsolenmenü, Fehlerbehandlung \\
    W3 & Sprint 1 & Sprint Review: JSON-Feedback, Englisch, Architektur \\
    W3 & Sprint 1 & Retrospektive: Erkenntnisse zu Story-Point-Schätzung und SRP \\
    W3 & Sprint 2 & Planning: 18 SP, JSON/English/Refactoring im Sprint-Ziel \\
    W3--W4 & Sprint 2 & Entwicklung: JsonExporter, Refactoring, JUnit-Tests \\
    W4 & Sprint 2 & Sprint Review: Abnahme durch Product Owner \\
    W4 & Sprint 2 & Retrospektive + Dokumentation begonnen \\
    W5 & Doku & Fertigstellung Dokumentation, Abgabe \\
    \bottomrule
  \end{tabularx}
\end{table}

\section{Konkrete Probleme während der Entwicklung}

\paragraph{Problem 1: Gson und Maven-Dependency (Sprint 2, Tag 1):}
Die Einbindung der Gson-Bibliothek via Maven führte zunächst zu einem
\texttt{ClassNotFoundException} beim ersten Testlauf. Ursache war eine
fehlende \texttt{scope}-Deklaration in der \texttt{pom.xml} und ein veraltetes
lokales Maven-Repository. Lösung: \texttt{mvn dependency:purge-local-repository}
sowie explizite \texttt{compile}-Scope-Angabe.

\paragraph{Problem 2: NullPointerException bei unbekannter ID (Sprint 1, Tag 3):}
Die \texttt{findById}-Methode in \texttt{Warehouse} gab \texttt{null} zurück,
wenn die ID nicht gefunden wurde. In \texttt{ConsoleUI} wurde das Ergebnis
direkt dereferenziert, was zu einem Absturz führte. Lösung: Umstellung auf
\texttt{Optional<Item>} und explizite Fehlerbehandlung in \texttt{ConsoleUI}.

\paragraph{Problem 3: Doppelte Verantwortlichkeit in \texttt{Main} (Sprint 1, Sprint Review):}
Die \texttt{Main}-Klasse enthielt anfangs sowohl die Menüsteuerung als auch
Teile der Geschäftslogik. Dies war der direkteste Verstoß gegen das
Single Responsibility Principle und wurde erst im Sprint Review explizit angemahnt.
Das Refactoring in Sprint~2 (Extraktion von \texttt{WarehouseService} und
\texttt{ConsoleUI}) behob dieses Problem vollständig.

\paragraph{Problem 4: Zeichencodierung in JSON (Sprint 2, Tag 3):}
Sonderzeichen in Artikelnamen (z.\,B. Anführungszeichen im Modellnamen
\textit{Dell UltraSharp 27'' Monitor}) führten ohne Escaping zu invaliden
JSON-Ausgaben. Gson behandelt das intern korrekt; das Problem trat allerdings
auf, als die Gruppe versuchte, selbst einen JSON-String zu konkatenieren.
Die vollständige Umstellung auf Gson's Serialisierungsmethoden behob das Problem.

%% ============================================================
%%  KAPITEL 20: FAZIT
%% ============================================================
\chapter{Fazit}
\label{chap:fazit}

\section{Zusammenfassung der Ergebnisse}

Das Praktikum \textit{Vom Wasserfall zum Sprint} hat sein didaktisches Ziel in vollem
Umfang erreicht. Das entwickelte Elektroniklagerverwaltungssystem ist funktional vollständig,
technisch sauber strukturiert und entspricht nach zwei Scrum-Sprints exakt den
Anforderungen des Product Owners. Die Systemarchitektur folgt dem Prinzip der Separation
of Concerns, alle 10 JUnit-Tests laufen fehlerfrei durch, und der JSON-Export produziert
valide, internationalisierbare Ausgaben.

\section{Methodische Erkenntnisse}

Die zentrale methodische Erkenntnis lässt sich in einem Satz zusammenfassen:
\textbf{Beide Modelle sind richtig \textendash{} für unterschiedliche Kontexte.}

Das Wasserfallmodell lieferte ein vollständiges System pünktlich nach Plan. Die
Planungsdisziplin war wertvoll. Die Erfahrung zeigte aber auch: \textit{Pünktlich am
falschen Ziel ankommen} ist kein Projekterfolg.

Scrum lieferte anfänglich mehr Orientierungslosigkeit und mehr Diskussion. Die
Sprint Reviews erzeugten Mehrarbeit. Aber diese Mehrarbeit war produktiv: Sie führte
zu einem System, das der Auftraggeber tatsächlich nutzen kann und das für spätere
Erweiterungen offen bleibt.

\section{Ausblick}

Das entwickelte System bildet eine solide Grundlage für weitere Erweiterungen:
eine REST-API-Anbindung (die saubere \texttt{WarehouseService}-Schicht ermöglicht dies
ohne Architekturbruch), eine datenbankgestützte Persistenz (Austausch von
\texttt{Warehouse} gegen eine JPA-Repository-Implementierung) oder eine Web-Benutzeroberfläche.
Die gewählte Architektur macht genau diese Erweiterbarkeit möglich \textendash{} ein Ergebnis,
das der Scrum-Prozess durch konsequentes Refactoring herbeigeführt hat.

\begin{infobox}[Abschließende Bewertung des Scrum-Einsatzes]
  Für ein System wie das vorliegende Elektroniklagerverwaltungssystem \textendash{} mit einem
  echten (simulierten) Auftraggeber, sich konkretisierenden Anforderungen und dem Ziel
  der späteren Erweiterbarkeit \textendash{} war Scrum das eindeutig überlegene Entwicklungsmodell.
  Nicht weil Wasserfall schlecht ist, sondern weil die Projektcharakteristika explizit
  auf die Stärken agiler Entwicklung zugeschnitten waren.
\end{infobox}

%% ============================================================
%%  LITERATURVERZEICHNIS
%% ============================================================
\printbibliography[heading=bibintoc, title={Literaturverzeichnis}]

\end{document}
